\chapter{Formal verification}\label{sec:formal}

The assembled Palomar submission
\url{https://github.com/pmaronow/fluctuations-palomar} exposes five
claims through \path{Solution.lean}, selected by \path{comparator.json}; its
independent \path{Challenge.lean} contains those five claims. The original
combined interface passed the recorded checks on 19 October 2026
(20 October UTC), for proof-source commit
\path{5551e62264a270fd5dcd396f36389eaf77eddebc}.
The default build succeeded using unchanged compiled objects, and a
strict axiom audit accepted all five exports with only
\path{propext}, \path{Quot.sound}, and \path{Classical.choice}.
A fresh protected Comparator run compiled the independent Challenge,
compared the five statements and their complete reachable helper
dependencies, and replayed the Solution export through Lean's kernel,
NanoDa, and con-ron. All three accepted it. The records are published at
\url{https://github.com/pmaronow/fluctuations-palomar/tree/351cc102439ff1d741d240391dfb62258605b4d2/verification}.
In that interface, the export for \cref{cor:mean-bias} contained the
lower-bound predicate alone. The current export also concludes convergence
of the finite-volume counting pressures to the defined limiting pressure.
Only \path{Challenge.lean} and \path{Solution.lean} change in the Lean
sources; the other four compared statements, the retained proof
development, and the four build/comparison controls are unchanged.
A fresh build of the strengthened interface at proof-source commit
\path{2fd7036fa3435aaacef2801c35ae95082d3f6806} completed successfully
on 20 October 2026, using Lean \texttt{v4.35.0-rc2} and Mathlib at
\path{065356127b1dc0016f66b7283ce0ce2c4055aa55}.
A strict audit using Lake's generated import manifest, with
\texttt{autoImplicit=false} and \texttt{warningAsError=true}, accepted
all five exports with only \path{propext}, \path{Quot.sound}, and
\path{Classical.choice}. The revised interface and these checks are recorded in
\url{https://github.com/pmaronow/fluctuations-palomar/blob/main/verification/corollary-1-7-convergence.md}.
The historical Comparator and three-kernel receipts remain scoped to their
named source snapshot. The source-package checks below, including the
historical 47-claim upper-bound check, retain their original scope.

In this section, we describe formal proofs of the four results
\cref{thm:barrier-730,b:cor:subcritical,f:main,cor:mean-bias} of the lower bound
and the fluctuation part in the Lean~4 proof assistant \cite{lean4}, built on the Mathlib library \cite{mathlib}. No argument elsewhere in the paper depends on this section. In \cref{sec:formal-oai}, we recall what OpenAI's public formalization proves. In \cref{sec:formal-statements}, we state the four formal theorems, explain how they correspond to the statements of this paper, and describe how they were checked. \Cref{sec:formal-barrier} compares the formal proof of \cref{thm:barrier-730} with the proof in this paper, and \cref{sec:formal-window} describes the mathematics that the formal proof of \cref{f:main} had to supply beyond OpenAI's formalization. \Cref{sec:formal-upper} describes the upper-bound formalization, and
\cref{sec:formal-scope} lists what is not formally verified.

\section{OpenAI's formalization}\label{sec:formal-oai}

OpenAI's public repository \cite{oai-math-lean} contains Lean formalizations accompanying \cite{oai-barriers,oai-fluctuations}. For the barrier, the theorem \path{OAI.SK.main}, in the file \path{lean/OAI/Probability/SKBarriers/Main.lean}, states that for every $\beta>1$ two quantities tend to one: the Gibbs mass, averaged over the disorder, of the initial states whose total variation distance from $\pi^J$ at continuous time $\exp(n^{1/10000})$ is greater than $1/4$, and the same mass with the distance after $\lfloor\exp(n^{1/10000})\rfloor$ discrete attempts. This is the formal counterpart of \cite[Theorem~1.1]{oai-barriers}; the repository's scope document \path{lean/docs/227.md} describes it.

For the fluctuations, the theorem \path{main_of_positiveLaplaceWindow}, at line~294 of \path{lean/OAI/Probability/SKFluctuations/EventuallyNatPos.lean}, has the hypothesis \path{hw : PositiveLaplaceWindow} and the conclusion \path{MainFluctuationStatement}. Both propositions are defined in the file \path{Spin.lean} of the same directory, at lines~453 and~531. The proposition \path{PositiveLaplaceWindow} is the statement of \cref{f:window}, for every $\beta>1$ and every $0<m<1/300$, and \path{MainFluctuationStatement} is the conjunction of the three assertions of \cref{f:main}, for every $\beta>1$ and every $\eta>0$. \path{MainFluctuationStatement} uses the log partition function before completion, normalized as $\log(2^{-n}\sum_xe^{\beta H_n(x)})=F_n-n\log2$, which has the same variance and the same centered fluctuations as $F_n$. \path{PositiveLaplaceWindow} uses its Gaussian completion $F_n^c-n\log2$, whose logarithmic moment generating function after centering is $K_n$. OpenAI's formalization proves the deduction of \cref{f:main} from \cref{f:window} given in \cite[Sections~7.5 and~7.6]{oai-fluctuations}, of which \cref{sec:fl-logconcave,sec:fl-endgame} give a variant, but not \cref{f:window} itself. Aronow and Lopatto inspected this directory and found no proof of the hypothesis there \cite[Section~1.6, footnote~1]{hexagon-bias}. In the paper \cite{oai-fluctuations}, by contrast, \cref{f:window} is proved, as summarized in \cref{sec:fl-route}.

\section{Statements of the formal results}\label{sec:formal-statements}

We originally formalized the four results in two source packages, one for \cref{thm:barrier-730,b:cor:subcritical} and one for \cref{f:main,cor:mean-bias}. Both use Lean version \texttt{v4.35.0-rc2} and Mathlib at commit \path{065356127b1dc0016f66b7283ce0ce2c4055aa55}. Both include OpenAI's Lean sources at commit \path{fd4aeeb2ee4fc729c18d98444fed42fd0529eeeb} as a dependency, adapted to the module system of this version of Lean. The second package also depends on the first, which contains the formal proof of the variance--bias inequality.

\paragraph{The barrier theorem.} The theorem \path{SKBarriers2.mainStatementSevenThirty} has type \path{OAI.SK.Strengthened.MainStatementAt (7 / 30)}. This proposition states that for every $\beta>1$ there exists $A>0$ with the following property: for every $\tau\in(0,1/2)$, the two quantities
\[
 \int\sum_{x\in\Sigma_n}\pi^J(x)\,\mathbf 1_{\{d^J(x,\exp(An^{7/30}))>\tau\}}\,\PP(\dd J),\qquad
 \int\sum_{x\in\Sigma_n}\pi^J(x)\,\mathbf 1_{\{d^J_{\rm disc}(x,\lfloor\exp(An^{7/30})\rfloor)>\tau\}}\,\PP(\dd J)
\]
tend to one as $n\to\infty$. These are the two probabilities in \cref{thm:barrier-730}, written as in \eqref{eq:common-joint}, and $A$ plays the role of $A_\beta$. As in \cref{thm:barrier-730}, the coefficient is chosen before the threshold $\tau$, and the same coefficient serves both clocks. The disorder law, the Hamiltonian, the Gibbs measure, the two heat-bath kernels, and the total variation distances are those of OpenAI's development. They agree with the definitions in \cref{sec:model}: the partition function $Z_n$ is a sum over $\Sigma_n$ with counting measure, one discrete step is an update attempt at a uniformly chosen site, an attempt that leaves the configuration unchanged is counted, and the continuous-time kernel is given by \eqref{eq:model-poisson}. Both displayed quantities are instances of a definition in the development that takes the coefficient, the exponent of $n$, and the threshold as arguments. The development proves that the instance with coefficient $1$, exponent $1/10000$, and threshold $1/4$ is the corresponding quantity of OpenAI's theorem. The theorem \path{SKBarriers2.subcriticalStatementSevenThirty}, of type \path{OAI.SK.Strengthened.SubcriticalStatementBelow (7 / 30)}, is \cref{b:cor:subcritical}: the same conclusion with coefficient one, at the time horizons $\exp(n^\kappa)$ and $\lfloor\exp(n^\kappa)\rfloor$, for every $0<\kappa<7/30$. The two theorems have no hypotheses.

\paragraph{The fluctuation theorem.} The fluctuation theorem is the declaration
\path{OAI.SKFluctuations.LaplaceWindowFormalization.main_fluctuation_statement_unconditional_original_q}
in the module \path{SKProof.OriginalQUnconditionalMain}. Its type is \path{OAI.SKFluctuations.MainFluctuationStatement}, the definition of OpenAI's development described in \cref{sec:formal-oai}, and it has no hypotheses. The identifier is only a locator; its parts are not used as terminology in this paper. In \path{MainFluctuationStatement}, the bounds on the standard deviation and on the variance come with separate thresholds in $n$, and their maximum serves as the threshold $n_0(\beta,\eta)$ of \cref{f:main}. The proof of the theorem establishes \path{PositiveLaplaceWindow} and then applies OpenAI's theorem \path{main_of_positiveLaplaceWindow}. In other words, our formal proof of \cref{f:main} consists of a formal proof of \cref{f:window} together with OpenAI's formal proof of the deduction.

\paragraph{The mean-bias corollary.} The compared export
\path{LowerUpperBounds.corollary_1_7} now has two conjuncts. The first
asserts, for every $\beta>1$, convergence of
\path{MeanBias.countingPressure beta} to
\path{MeanBias.limitingPressure beta}. Here and below these names are in
\path{OAI.SKFluctuations}. The limiting-pressure definition uses
\path{limUnder atTop}; the convergence is an explicit conclusion of the
export, not part of that definition. The second conjunct is the unchanged
proposition \path{OAI.SKFluctuations.MeanBias.LowerBoundStatement}.

The lower-bound conjunct is supplied by the library declaration
\path{OAI.SKFluctuations.LaplaceWindowFormalization.mean_bias_lower_bound_unconditional},
in the module \path{SKProof.MeanBiasLowerBound}, whose type remains
\path{OAI.SKFluctuations.MeanBias.LowerBoundStatement}. Together with the
convergence conjunct, it states that for every $\beta>1$ and every
$\eps>0$, there exist $c>0$ and $n_0$ such that
\[
 c\,n^{-5/6-\eps}\le\lim_{m\to\infty}p_m(\beta)-p_n(\beta)\qquad(n\ge n_0),
\]
where we write $p_n(\beta)=n^{-1}\E\log\sum_{x\in\Sigma_n}e^{\beta H_n(x)}$, with OpenAI's disorder law and Hamiltonian, which are those of \cref{sec:model}. Since $P_\beta=\lim_mp_m(\beta)$ by definition, this is \cref{cor:mean-bias}. To prove the convergence conjunct, \path{Solution.lean} identifies \path{MeanBias.countingPressure beta} with \path{Paper.finiteSKFreeEnergy beta 0} using the pressure identities in \path{SKProof.MeanBiasBridge} and \path{QuantitativePaper.Model}. It then applies \path{Paper.physicalParisiFormula}, in \path{Paper.ParisiPDEFormula}, which proves convergence to the infimum of the Parisi functional for every positive $\beta$. The bridge's theorem \path{MeanBias.limitingPressure_eq_parisiPDEValue} identifies this value with \path{limitingPressure beta}, using that convergence. Thus the compared export includes existence of the actual pressure limit as well as the lower bound; the underlying library lower-bound predicate is unchanged.

To prove this bound, we follow the proof of \cref{cor:mean-bias} in \cref{sec:bias-proof}. We combine our formal proof of \cref{f:main} with the theorem \path{QuantitativePaper.short_variance_mean} of the Lean development \cite{hexagon-bias-lean} accompanying \cite{hexagon-bias}, which is a formal version of \cref{prop:bias-variance-deficit}. That theorem concerns the variable $F_n/\beta$: for every $\beta_0>1$, there exists $c>0$ such that
\[
 c\Bigl(\frac{\Var(F_n/\beta)+1/2}{n}\Bigr)^{5/4}\le\frac{\Delta_{n,\beta}}{\beta}
\]
for every $n\ge1$ and every $\beta\ge\beta_0$. Since $\Var(F_n/\beta)+1/2=\beta^{-2}(\Var(F_n)+\beta^2/2)$, at a fixed $\beta$ this is \eqref{eq:mean-bias-variance-deficit} with $c_\beta=c\beta^{-3/2}$. The module \path{SKProof.MeanBiasBridge} proves the identities that connect the two developments: their disorder laws and Hamiltonians coincide, their expected pressures agree, and OpenAI's log partition function is $F_n-n\log2$, so that its variance is $\Var(F_n)$. As in \cref{sec:bias-proof}, the lower bound on the variance in \cref{f:main}, with $\eta=4\eps/5$, then gives the exponent $-5/6-\eps$.

\paragraph{Axioms and the independent check.} The four original source-package theorems use only the three standard axioms \path{propext}, \path{Classical.choice}, and \path{Quot.sound} of Lean's logic. Since they also have no hypotheses, every result that their proofs use is proved in Lean, in the two packages or in their dependencies. The upper bounds on the mean bias and on the variance in \cref{inp:hb-rate,inp:qp-rate} and the variance--bias inequality enter the formal proofs as Lean theorems. The first and the last of these come from \cite{hexagon-bias-lean}. The strengthened combined interface passed the strict five-export axiom audit described at the start of this section.

In their source packages, the four theorems also passed an independent check, which has three parts. First, we wrote out their statements a second time, with every definition that they involve, in separate files that import only Mathlib and none of our proofs. For the barrier, these are the definitions of the disorder law, the Hamiltonian, the Gibbs measure, the heat-bath kernels, the total variation distances, and the two quantities displayed above. For the fluctuations and the mean bias, they are the definitions on which \path{MainFluctuationStatement} rests, copied from OpenAI's sources at the commit cited above, and the definitions of the pressures $p_n(\beta)$ and of their limit. A comparison confirmed that each proved theorem has exactly the statement written out in this way. Second, the same comparison confirmed that the proofs use no axioms other than the three above. Third, the statements and proofs of the theorems were exported from the compiled packages, with every declaration on which they depend, including the proofs of all lemmas that they use. This export was checked again by three type checkers for Lean's logic: Lean's own kernel and the external checkers nanoda and con-ron. All three accepted the barrier statements and, in a separate comparison, both the fluctuation theorem and the mean-bias corollary. Declarations of the packages that these proofs do not use are not part of the check.

\section{The formal and the written barrier proofs}\label{sec:formal-barrier}

The formal proof of \cref{thm:barrier-730} does not follow the proof in this paper in every step. The paper proves the local three-replica expansion of \cref{b:prop:os-expansion} with a remainder $O(h^4k^2)$, where $h$ is the width parameter and $k$ the relative imbalance of \eqref{b:eq:os-parameters}, by estimates for smooth limiting flows. The formal proof instead works with finite chains of Gaussian integrations, using the hierarchical recursion of \cite[Section~2]{oai-barriers} and its accompanying Lean development \cite{oai-math-lean}. It proves the upper bound contained in \eqref{b:eq:finiteE-polynomial} for these chains, at the ratio $\vartheta=1/2$ used in the dynamical argument and on a smaller range of the relative imbalance $k$. Its remainder is of order $h^{7/2}k^2+h^3|k|^3$ instead of $h^4k^2$, up to an additional error that tends to zero as the length of the chains tends to infinity. Only this upper bound enters the proof of the uniform cubic deficit in \cref{prop:cub-band}, where the remainder must be small compared with $h^3k^2$. The weaker remainder has this property when $h$ and $|k|$ are small. Either estimate suffices for \cref{thm:barrier-730}; they lead to different constants, not to a different exponent.

In the final counting step, the formal proof and \cref{sec:barrier-proof} use the same kind of count. Only the first reference configuration of each block is counted, and the sign obstruction follows, as in \cref{prop:bp-obstruction}, once a count of the form \eqref{b:dyn:prefix-count} is less than one. The formal proof passes through finite grids of threshold values, while \cref{sec:barrier-proof} uses the uniform law of the thresholds directly. We conclude that the formal proof certifies \cref{thm:barrier-730} and \cref{b:cor:subcritical}, not every intermediate estimate in the paper's proof of them.

\section{The formal proof of the two-sided bound}\label{sec:formal-window}

Our formal proof of \cref{f:window} follows the route of \cite{oai-fluctuations} summarized in \cref{sec:fl-route}, and its mathematics is that of \cite{oai-fluctuations}. We do not claim a new proof of \cref{f:window} or of \cref{f:main}, and we do not claim to have found or repaired a gap in the argument of \cite{oai-fluctuations}. What the formalization supplies is a complete treatment of the particular variational problems to which \cite[Section~7]{oai-fluctuations} applies the general results of \cite[Sections~3--6]{oai-fluctuations}. A paper proof can handle these problems by checking that they satisfy the hypotheses of general statements. In the formal proof, each such check is carried out in full detail. We now describe these problems and the cases that arise, in the notation of \cref{sec:fl-route}.

\paragraph{The comparisons and their minimizers.} The upper bound in \cref{f:window} rests on the comparisons of \cite[Section~7.3]{oai-fluctuations}. For each $\eps\in[\eps_0,\beta^2)$ and $\theta\in[0,1]$, we minimize a functional of the form \cite[(6)]{oai-fluctuations} over overlap quantiles $r$ on the mass interval $[s,1]$. The lower radius is $r(s)=r_{\min}={\bar u}_*/(2M(s))$, and the low-field perturbation of strength $\theta$ acts on $[0,s)$. The general results of \cite{oai-fluctuations} are stated for the comparison families of \cite[Definition~6.1]{oai-fluctuations}, which allow a positive lower radius $r_{\min}$ and an omitted lower interval. All overlap moments in them refer to minimizers of the perturbed paths. In the formal proof, we establish their hypotheses and conclusions for the minimizers of these particular functionals and for the finitely many successive comparisons built from them. For $\theta=0$, and for the successive comparisons, whose unperturbed paths are nonnegative and nondecreasing, we do this throughout the strength interval $0<{\widehat\lambda}\le d^2$. For $\theta>0$, we use the two-sided linear bound only at the fixed final strength ${\widehat\lambda}=d^2$, uniformly in $\eps$, $\theta$, and the minimizer. We use neither a pressure at the unperturbed paths nor an integration in the strength from zero. This restriction is necessary. Let $e_{\widehat\lambda}$ be the function $e$ of \cref{sec:fl-route} at strength ${\widehat\lambda}$. Since $e_{\widehat\lambda}(r_{\min})={\widehat\lambda}$ for large $n$, the perturbed field path is nondecreasing at $s$ only if ${\widehat\lambda} r_{\min}\ge\theta c_Lr_0$, that is, only if ${\widehat\lambda}\ge\theta c_Ld^2/(2\eps)$.

In the formal proof, we use the strength profile $e_{\widehat\lambda}(r)={\widehat\lambda}\tilde f({\widehat\lambda} r/d^3)$, where $\tilde f(t)=1$ for $t\le1$ and $\tilde f(t)=(1+\log t)/t$ for $t>1$. This function has the monotonicity properties required of the profile $f$ in \cite[(6)]{oai-fluctuations}: its logarithmic slope $-t\tilde f'(t)/\tilde f(t)$ equals $0$ for $t\le1$ and $\log t/(1+\log t)$ for $t>1$. It is continuously differentiable but not smooth at $t=1$, and the parts of the argument that use the profile were carried out for this choice.

\paragraph{The three cases.} Consider a minimizer at strength ${\widehat\lambda}$ and a mass $x\ge A_{\widehat\lambda}$ at which the two-sided linear bound fails, while it holds above $2x$, as in the propagation argument of \cref{sec:fl-route}. There are three cases.
\begin{enumerate}
\item \emph{Comparable strength at a small scale}: $x\to0$ and $e_{\widehat\lambda}(x)/x^2$ is bounded below. This case is excluded by the linear limit of \cite[Theorem~5.1]{oai-fluctuations}, which has to be proved for minimizers with the positive lower radius $r_{\min}$. Terms involving $r_{\min}$ enter as multiples of ${\widehat\lambda} r_{\min}$, and they are negligible at the cubic scale. To see this, recall from \cite[(3) and (5)]{oai-fluctuations} that
\[
 \underline x=\frac{\sqrt nE}{d},\qquad {\bar u}_*=\frac1{nE},\qquad L=\min\Bigl\{{\bar u}_*,\frac{\underline x^3}{d^2}\Bigr\},
\]
and that the upper bound on the derivative of the inverse quantile is
\[
 M({\bar u})=n^\delta\Bigl[\frac{{\bar u}+{\bar u}_*}{L}+\frac{d^3}{E}({\bar u}+{\bar u}_*)^2\Bigr].
\]
Since $M(s)\ge n^\delta {\bar u}_*/L$ and $r_{\min}={\bar u}_*/(2M(s))$, and since $x\ge A_{\widehat\lambda}\ge n^\delta\underline x\ge\underline x$ for ${\widehat\lambda}\le d^2$, we have
\[
 {\widehat\lambda} r_{\min}\le\frac{{\widehat\lambda} L}{2n^\delta}\le\frac{d^2L}{2n^\delta}\le\frac{\underline x^3}{2n^\delta}\le\frac{x^3}{2n^\delta}.
\]
We use the one-site cavity estimate of \cite[Proposition~4.1]{oai-fluctuations} at the two scales $x$ and $\bar x=\max\{x,\sqrt{e_{\widehat\lambda}(x)}\}$, and $\bar x/x$ may tend to infinity; for example, $\bar x/x=n^{5m-\delta}$ when ${\widehat\lambda}=d^2$ and $x=A_{d^2}$. The identification of the limiting scalar and overlap distributions \cite[Lemma~5.4]{oai-fluctuations} and the first variations on gaps of the limiting support \cite[Lemma~5.6]{oai-fluctuations} are carried out for these minimizers and their admissible competitors.
\item \emph{Vanishing strength}: $e_{\widehat\lambda}(x)/x^2\to0$, either with $x\to0$ or with $x$ bounded below. The next comparison is built from the perturbed paths of this minimizer, as in \cite[Lemma~6.2]{oai-fluctuations}. The alternative described in \cref{sec:fl-route} is then followed for a fixed number $N=\lceil10/\delta\rceil$ of successive comparisons \cite[Lemmas~6.3--6.5 and Proposition~6.6]{oai-fluctuations}. At each stage, we exclude a violation with strength comparable to the scale as in the first or the third case.
\item \emph{Macroscopic scale}: $x$ and $e_{\widehat\lambda}(x)/x^2$ are bounded below, and $d^2\le\beta^2-c_*$ for a fixed $c_*>0$. We exclude this case using the convergence of the minimizers to the quantile function of the Parisi measure \cite[Proposition~5.9]{oai-fluctuations}, applied to the comparison at hand.
\end{enumerate}
The successive comparisons follow one of the two covariance alternatives of \cite[Definition~6.1 and Corollary~6.7]{oai-fluctuations}. If $\eps<\beta^2/4$, the perturbed matrix covariance is bounded below by a fixed positive constant, and they act on all masses. If $\eps\ge\beta^2/4$, they act only on $[s,1]$, with lower radius zero, and the matrix covariance has a jump of at least $\beta^2/4-d^2$ at $s$. We then choose the coefficient $c_0$ so that the total matrix covariance $Nc_0^2$ removed by the successive comparisons is less than half of a fixed positive constant that bounds this jump from below; their perturbed matrix paths therefore remain nondecreasing at $s$. For $\theta>0$, we establish the two-sided linear bound at the fixed final strength $d^2$ \cite[Lemma~7.2]{oai-fluctuations} and the bound $S_{\bar u}\le Cd^3/\eps$ on $[0,s)$ \cite[Lemma~7.3]{oai-fluctuations}, uniformly in $\eps$ and $\theta$. These two estimates enter the integration in $\theta$ of \cite[Lemma~7.4]{oai-fluctuations}. As in \cite{oai-fluctuations}, the constants in the two-sided linear bound, the depth $N$, and the coefficient $c_0$ are fixed before $n\to\infty$. The final threshold in $n$ is uniform in $\eps$, $\theta$, and the minimizer, and in ${\widehat\lambda}$ over the strengths at which each comparison is used.

\section{The upper bound}\label{sec:formal-upper}
\statusFormalUpper

\section{Scope and availability}\label{sec:formal-scope}

The combined submission selects the five final statements described above, including the strengthened mean-bias export checked by the fresh build and strict audit above; the historical source-package checks of \cref{sec:formal-statements} concern the four original lower-bound and fluctuation results described there. Intermediate results in their proofs, such as \cref{f:window}, are checked only as parts of these proofs. The quantitative and almost-sure forms in \cref{thm:barrier-strong} are not formally verified. These are the bounds on the probability of the exceptional disorders and on the Gibbs mass of the exceptional initial states, the lower bounds on the distance that tend to $1/2$, the extension to the Metropolis dynamics and to other single-site dynamics with bounded rates, the bound on the relaxation time, and the almost-sure statements for arbitrary couplings of the disorders across system sizes. The sharper local expansion of \cref{b:prop:os-expansion} is not formally verified either; see \cref{sec:formal-barrier}. The proposition \path{MainFluctuationStatement} concerns exponents only, and the limit law and the positive variance prefactor claimed in \cite{oai-limiting-law} are not part of it. The proposition \path{LowerBoundStatement} contains the loss $n^{-\eps}$ for an arbitrary $\eps>0$, and the bound \eqref{eq:mean-bias-prefactor-target}, which would follow from a positive variance prefactor, is not part of it.

\statusFormalUpperScope

OpenAI's Lean sources \cite{oai-math-lean} and the development
\cite{hexagon-bias-lean} are public. The combined sources, independent
\path{Challenge.lean}, and verification records are available in
\url{https://github.com/pmaronow/fluctuations-palomar}; the records
identify the checked proof-source revision and the scope of each check.
The original two source packages and the historical upper-bound package
are maintained in the private repository
\url{https://github.com/pmaronow/fluctuations-temp}. The latter was checked
at commit \path{352426b9007b6bbb31a3c8cfe6887caa0badf3fd}, in its
\path{glauber/lean} directory. The historical records remain available
from the authors.
